\documentclass[12pt,a4paper]{article}
\usepackage[utf8]{inputenc}
\usepackage[margin=0.75in]{geometry}
\usepackage{xcolor}
\usepackage{tcolorbox}
\usepackage{enumitem}
\usepackage{titlesec}
\usepackage{fancyhdr}
\usepackage{hyperref}

% Color definitions
\definecolor{slidecolor}{RGB}{52,152,219}
\definecolor{emphasiscolor}{RGB}{22,160,133}
\definecolor{quotecolor}{RGB}{249,249,249}
\definecolor{tipcolor}{RGB}{213,244,230}

% Section formatting
\titleformat{\section}
  {\normalfont\LARGE\bfseries\color{slidecolor}}{\thesection}{1em}{}
\titleformat{\subsection}
  {\normalfont\Large\bfseries\color{emphasiscolor}}{\thesubsection}{1em}{}

% Header and footer
\pagestyle{fancy}
\fancyhf{}
\fancyhead[C]{\textbf{Slide-by-Slide Delivery Guide - Conformal Mappings}}
\fancyfoot[C]{\thepage}
\renewcommand{\headrulewidth}{0.4pt}

% Custom boxes
\newtcolorbox{quotebox}{
  colback=quotecolor,
  colframe=slidecolor,
  boxrule=0pt,
  leftrule=5pt,
  arc=0pt,
  left=12pt,
  right=12pt,
  top=8pt,
  bottom=8pt
}

\newtcolorbox{tipbox}{
  colback=tipcolor,
  colframe=emphasiscolor,
  boxrule=1.5pt,
  arc=3pt,
  left=12pt,
  right=12pt,
  top=8pt,
  bottom=8pt
}

\begin{document}

\begin{center}
{\Huge \textbf{SLIDE-BY-SLIDE DELIVERY GUIDE}}\\[10pt]
{\Large Computation of Conformal Mappings Presentation}\\[8pt]
{\large \textbf{BTP Defense TODAY at 1:00 PM - 60\% Marks}}
\end{center}

\vspace{10pt}
\hrule
\vspace{20pt}

\section*{SLIDE 1: Title Slide}

\subsection*{What to Say:}
\begin{quotebox}
\large\textit{``Good afternoon professors. Today I'll present my work on the Computation of Conformal Mappings, specifically focusing on the Schwarz-Christoffel transformation. This is a powerful mathematical tool that bridges pure theory and practical engineering applications.''}
\end{quotebox}

\subsection*{Body Language:}
\begin{itemize}[leftmargin=*]
    \item Stand confidently
    \item Make eye contact with all professors
    \item Pause after greeting
\end{itemize}

\vspace{10pt}
\hrule
\vspace{10pt}

\section*{SLIDE 2: Why Conformal Mapping?}

\subsection*{What to Say:}
\begin{quotebox}
\large\textit{``Let me start with a fundamental question: Why do we need conformal mapping? In engineering and physics, we often face problems with irregular, complex boundaries---think of airflow around an oddly-shaped wing or heat distribution in an irregular metal plate.}

\textit{Traditional analytical methods only work in simple geometries like circles or rectangles. Conformal mapping is the \textbf{bridge} that transforms these complex domains into simple ones while preserving the governing equations. The diagram on the right shows this visually---we solve the problem in the simple Z-plane and map the solution back.''}
\end{quotebox}

\subsection*{Key Points to Emphasize:}
\begin{itemize}[leftmargin=*]
    \item Use the diagram---point to it while explaining
    \item Say ``This is \textbf{THE} solution to an otherwise intractable problem''
    \item Pause after ``preserving governing equations''
    \item Connect to real-world problems (airflow, heat transfer)
\end{itemize}

\vspace{10pt}
\hrule
\vspace{10pt}

\section*{SLIDE 3: The Riemann Mapping Theorem}

\subsection*{What to Say:}
\begin{quotebox}
\textit{``The theoretical foundation is the Riemann Mapping Theorem, proved by Riemann in 1851 and rigorously established by Schwarz and Hilbert.}

\textit{The theorem states: \textbf{Any simply connected domain---no matter how irregular---can be conformally mapped to the unit disk.} This is profound because it guarantees existence. However, and this is crucial, it doesn't tell us \textbf{how} to construct the mapping. That's where Schwarz-Christoffel comes in.''}
\end{quotebox}

\subsection*{Key Points to Emphasize:}
\begin{itemize}[leftmargin=*]
    \item Point to the diagram showing the irregular domain mapping to the disk
    \item Emphasize ``existence vs.\ construction''---this sets up the need for SC transformation
    \item Make clear the historical significance (1851, later proved rigorously)
    \item Connect: ``The theorem says it CAN be done, SC shows us HOW''
\end{itemize}

\vspace{10pt}
\hrule
\vspace{10pt}

\section*{SLIDE 4: Foundations of Complex Analysis}

\subsection*{What to Say:}
\begin{quotebox}
\textit{``To understand conformal mappings, we need three key concepts from complex analysis:}

\textit{First, \textbf{analytic functions}---functions where the derivative exists everywhere in a region. Second, the \textbf{Cauchy-Riemann equations}, which are the necessary and sufficient conditions for analyticity. And third, \textbf{harmonic functions}---the real and imaginary parts satisfy Laplace's equation.}

\textit{This connection to Laplace's equation is why conformal mapping works for problems in electrostatics, fluid flow, and heat transfer---they're all governed by Laplace's equation.''}
\end{quotebox}

\subsection*{Key Points to Emphasize:}
\begin{itemize}[leftmargin=*]
    \item Don't rush this slide---it's foundational
    \item Point to the visual showing orthogonal grids
    \item Emphasize the connection: analytic $\rightarrow$ harmonic $\rightarrow$ Laplace's equation
    \item If time permits, mention: ``The orthogonality shown here is a direct consequence of the Cauchy-Riemann equations''
\end{itemize}

\vspace{10pt}
\hrule
\vspace{10pt}

\section*{SLIDE 5: Conformal Mapping Definition}

\subsection*{What to Say:}
\begin{quotebox}
\textit{``What makes a mapping conformal? It must preserve angles---both magnitude and orientation. The sufficient condition is elegant: if $f(z)$ is analytic and its derivative is non-zero, then $f$ is conformal.}

\textit{Look at the diagram: two curves intersect at angle $\alpha$ in the z-plane. After the conformal map $f(z)$, they still intersect at the same angle $\alpha$ in the w-plane. This angle preservation is what allows physical laws to remain unchanged under the transformation.''}
\end{quotebox}

\subsection*{Key Points to Emphasize:}
\begin{itemize}[leftmargin=*]
    \item Visual aid: Trace the angle with your hand on the diagram
    \item Emphasize the condition: $f'(z) \neq 0$ (local univalence)
    \item Connect to physics: ``Angle preservation means the form of PDEs is preserved''
    \item Mention the key property box on orthogonality if there's time
\end{itemize}

\vspace{10pt}
\hrule
\vspace{10pt}

\section*{SLIDE 6: Bilinear (Möbius) Transformation}

\subsection*{What to Say:}
\begin{quotebox}
\textit{``Before tackling polygons, let's look at the simplest conformal map: the Bilinear or Möbius transformation, $w = \frac{az+b}{cz+d}$.}

\textit{Its superpower is that it maps circles and lines to circles and lines. This makes it perfect for preliminary transformations---for instance, mapping the upper half-plane to the unit disk. Any Bilinear transformation can be decomposed into translation, inversion, and rotation.''}
\end{quotebox}

\subsection*{Key Points to Emphasize:}
\begin{itemize}[leftmargin=*]
    \item Keep it brief---this is a stepping stone to SC
    \item Point to the diagram showing circle/line preservation
    \item Mention the condition: $ad - bc \neq 0$ (ensures invertibility)
    \item Connect: ``We'll use this as a building block in more complex mappings''
\end{itemize}

\vspace{10pt}
\hrule
\vspace{10pt}

\section*{SLIDE 7: Schwarz-Christoffel Transformation - Purpose}

\subsection*{What to Say:}
\begin{quotebox}
\textit{``Now we arrive at the main contribution: the Schwarz-Christoffel transformation. It maps the upper half-plane onto the interior of \textbf{any} simple polygon.}

\textit{The formula is shown here. The key insight is in the product term: each factor $(z - x_k)^{-\alpha_k}$ controls what happens at a vertex. The points $x_k$ on the real axis map to the polygon vertices, and $\alpha_k$ are the exterior angles.''}
\end{quotebox}

\subsection*{Key Points to Emphasize:}
\begin{itemize}[leftmargin=*]
    \item \textbf{Slow down here---this is the core}
    \item Write or point to the formula: $\frac{dw}{dz} = A \prod (z - x_k)^{-\alpha_k}$
    \item Explain each component:
    \begin{itemize}
        \item $x_k$: pre-images on real axis
        \item $\alpha_k$: exterior angle parameters
        \item $A, B$: normalization constants
    \end{itemize}
    \item Emphasize: ``This explicit formula is what makes SC so powerful''
\end{itemize}

\vspace{10pt}
\hrule
\vspace{10pt}

\section*{SLIDE 8: Schwarz-Christoffel Derivation Concept}

\subsection*{What to Say:}
\begin{quotebox}
\textit{``How does this formula work? The elegant idea is to control the \textbf{argument} (angle) of $\frac{dw}{dz}$ as $z$ moves along the real axis.}

\textit{As $z$ crosses a point $x_k$, the factor $(z - x_k)^{-\alpha_k}$ causes the argument to change by exactly $\alpha_k\pi$---the exterior angle needed at that vertex. The diagram shows this: as $z$ moves along the real axis in the upper half-plane, $w$ traces the polygon boundary, turning by the correct angle at each vertex.''}
\end{quotebox}

\subsection*{Key Points to Emphasize:}
\begin{itemize}[leftmargin=*]
    \item \textbf{This is the `aha' moment---make it clear!}
    \item Use the diagram extensively---show the correspondence
    \item Explain the argument change: ``When $z$ crosses $x_k$ from left to right, $\arg(z - x_k)$ changes by $-\pi$, so $\arg((z - x_k)^{-\alpha_k})$ changes by $\alpha_k\pi$''
    \item Connect boundary mapping: ``Real axis $\rightarrow$ polygon boundary''
    \item Emphasize: ``This is the genius of Schwarz and Christoffel---encoding geometry in the derivative''
\end{itemize}

\vspace{10pt}
\hrule
\vspace{10pt}

\section*{SLIDE 9: Schwarz-Christoffel Construction Steps}

\subsection*{What to Say:}
\begin{quotebox}
\textit{``To construct an SC mapping, we follow four steps:}

\textit{1. Identify where on the real axis each vertex maps to---these are the $x_k$ points}

\textit{2. Determine the exterior angles from the polygon's interior angles using $\alpha_k = 1 - \theta_k/\pi$}

\textit{3. Construct the product ensuring each factor contributes the correct turn}

\textit{4. Integrate to get $w$, where $A$ and $B$ control scale, rotation, and position.''}
\end{quotebox}

\subsection*{Key Points to Emphasize:}
\begin{itemize}[leftmargin=*]
    \item Go through each step methodically
    \item Point to the diagram showing the UHP and polygon
    \item If asked about the integral: ``The integral often involves elliptic functions, which is why numerical methods are typically used in practice''
    \item Mention the constraint: $\sum\alpha_k = 2$ (polygon closure condition)
\end{itemize}

\vspace{10pt}
\hrule
\vspace{10pt}

\section*{SLIDE 10: Rectangle Mapping Example}

\subsection*{What to Say:}
\begin{quotebox}
\textit{``Let's see a concrete example: mapping to a rectangle. All four interior angles are $\pi/2$, so all exterior angle parameters $\alpha_k = 1 - (\pi/2)/\pi = 1/2$.}

\textit{The differential becomes $\frac{dw}{dz} = \frac{A}{\sqrt{(z-x_1)(z-x_2)(z-x_3)(z-x_4)}}$. This integral is an \textbf{elliptic integral}---it cannot be expressed in elementary functions. This shows that even for a simple rectangle, we need special functions or numerical methods.''}
\end{quotebox}

\subsection*{Key Points to Emphasize:}
\begin{itemize}[leftmargin=*]
    \item Point out the visual: Show how the UHP maps to the rectangle
    \item Highlight the four vertices and their pre-images
    \item Emphasize: ``Even simple polygons lead to complex integrals''
    \item If time permits: ``The ratio of rectangle sides depends on the ratio of complete elliptic integrals''
\end{itemize}

\vspace{10pt}
\hrule
\vspace{10pt}

\section*{SLIDE 11: Schwarz-Christoffel More Examples}

\subsection*{What to Say:}
\begin{quotebox}
\textit{``Slide 11 shows three more examples demonstrating the versatility of SC mappings:}

\textit{\textbf{Example 1}: The unit disk---while typically handled by Bilinear transformations, it shows the connection between different conformal maps.}

\textit{\textbf{Example 2}: A semi-infinite strip where the integral simplifies to $\cosh^{-1}(z)$, giving us a closed-form solution.}

\textit{\textbf{Example 3}: A triangular region with specific angles $\pi/4, \pi/4, \pi/2$, showing how different angle combinations affect the formula.''}
\end{quotebox}

\subsection*{Key Points to Emphasize:}
\begin{itemize}[leftmargin=*]
    \item Don't linger---these are for completeness
    \item Mention that some cases yield elementary functions (strip), others don't (general triangle)
    \item If asked: ``The semi-infinite strip is important for channel flow problems''
\end{itemize}

\vspace{10pt}
\hrule
\vspace{10pt}

\section*{SLIDE 12: Real-World Applications}

\subsection*{What to Say:}
\begin{quotebox}
\textit{``Why does this matter? Conformal mapping has revolutionized multiple fields:}

\textit{In \textbf{fluid dynamics}, the Joukowsky transformation maps circles to airfoil shapes, enabling the calculation of lift. In \textbf{electrostatics}, it solves for field lines around complex conductor shapes. In \textbf{heat transfer}, it analyzes temperature distributions in irregular domains. And in modern \textbf{computational engineering}, it generates high-quality grids for numerical simulations.''}
\end{quotebox}

\subsection*{Key Points to Emphasize:}
\begin{itemize}[leftmargin=*]
    \item Be enthusiastic here---show real-world impact!
    \item Point to each application icon in the diagram
    \item Connect to the central image showing the five applications
    \item If time permits, elaborate on one: ``In aerodynamics, conformal mapping was crucial before CFD---it gave analytical predictions of lift forces''
\end{itemize}

\vspace{10pt}
\hrule
\vspace{10pt}

\section*{SLIDE 13: Conclusion and Future Directions}

\subsection*{What to Say:}
\begin{quotebox}
\textit{``To conclude, conformal mapping---anchored by the Riemann Mapping Theorem and made explicit through the Schwarz-Christoffel formula---provides a powerful bridge between theory and application. It transforms intractable boundary value problems into solvable ones.}

\textit{Looking forward, future work includes:}
\begin{itemize}
    \item \textit{Validation through simulations and numerical examples}
    \item \textit{Developing faster numerical algorithms for the parameter problem}
    \item \textit{Advanced techniques for multiply-connected domains}
    \item \textit{Applications to modern boundary-value problems in engineering}
\end{itemize}

\textit{The interplay between theory and computation continues to drive innovation in this field. Thank you for your attention. I'm happy to answer questions.''}
\end{quotebox}

\subsection*{Key Points to Emphasize:}
\begin{itemize}[leftmargin=*]
    \item Summarize the three key boxes: Theoretical Foundation, Key Transformations, Conformal Mapping
    \item Highlight the future scope bullet points
    \item End with confidence and openness to questions
    \item Maintain eye contact during the closing
\end{itemize}

\vspace{10pt}
\hrule
\vspace{10pt}

\section*{GENERAL DELIVERY TIPS}

\begin{tipbox}
\subsection*{Voice \& Pace:}
\begin{itemize}[leftmargin=*]
    \item Speak at 120-130 words per minute (moderate pace)
    \item Pause after key statements (2-3 seconds)
    \item Vary your tone---don't be monotone
    \item Slow down for complex equations
\end{itemize}
\end{tipbox}

\begin{tipbox}
\subsection*{Body Language:}
\begin{itemize}[leftmargin=*]
    \item Stand tall, don't lean on the podium
    \item Use hand gestures to emphasize points
    \item Make eye contact with all professors (rotate your gaze)
    \item Point to diagrams and equations as you discuss them
\end{itemize}
\end{tipbox}

\begin{tipbox}
\subsection*{Engagement Strategies:}
\begin{itemize}[leftmargin=*]
    \item Use the pointer/cursor to guide attention
    \item Ask rhetorical questions: ``Why is this important?''
    \item Connect slides: ``Building on what we just saw...''
    \item Use transitional phrases: ``This leads us to...'', ``The key insight is...''
\end{itemize}
\end{tipbox}

\begin{tipbox}
\subsection*{Mathematical Rigor:}
\begin{itemize}[leftmargin=*]
    \item Don't skip over equations---explain at least the main SC formula in detail
    \item If asked to derive something, confidently walk through it
    \item Be prepared to explain any symbol or term on your slides
    \item Have alternate explanations ready for complex concepts
\end{itemize}
\end{tipbox}

\vspace{10pt}
\hrule
\vspace{10pt}

\section*{SLIDE TRANSITIONS}

\begin{itemize}[leftmargin=*]
    \item \textbf{Slide 1 $\rightarrow$ 2}: ``Let me begin by motivating why this topic matters...''
    \item \textbf{Slide 2 $\rightarrow$ 3}: ``The theoretical foundation that makes all this possible is...''
    \item \textbf{Slide 3 $\rightarrow$ 4}: ``To understand how this works, we need some complex analysis...''
    \item \textbf{Slide 4 $\rightarrow$ 5}: ``These concepts lead us to the definition of conformal mapping...''
    \item \textbf{Slide 5 $\rightarrow$ 6}: ``A simple but powerful example of a conformal map is...''
    \item \textbf{Slide 6 $\rightarrow$ 7}: ``Now we're ready for the main result...''
    \item \textbf{Slide 7 $\rightarrow$ 8}: ``But how does this formula actually work? Let's derive the concept...''
    \item \textbf{Slide 8 $\rightarrow$ 9}: ``Let me break down the construction process...''
    \item \textbf{Slide 9 $\rightarrow$ 10}: ``Let's make this concrete with an example...''
    \item \textbf{Slide 10 $\rightarrow$ 11}: ``Here are additional examples showing the range of applications...''
    \item \textbf{Slide 11 $\rightarrow$ 12}: ``All of this theory translates into powerful real-world applications...''
    \item \textbf{Slide 12 $\rightarrow$ 13}: ``In conclusion...''
\end{itemize}

\vspace{10pt}
\hrule
\vspace{10pt}

\section*{TIME MANAGEMENT (for $\sim$15-minute presentation)}

\begin{center}
\begin{tabular}{|l|c|}
\hline
\textbf{Slide} & \textbf{Time} \\
\hline
Slide 1 & 1 minute \\
Slide 2 & 2 minutes \\
Slide 3 & 1.5 minutes \\
Slide 4 & 1.5 minutes \\
Slide 5 & 1.5 minutes \\
Slide 6 & 1 minute \\
Slide 7 & 2 minutes (KEY SLIDE) \\
Slide 8 & 2 minutes (KEY SLIDE) \\
Slide 9 & 1.5 minutes \\
Slide 10 & 1.5 minutes \\
Slide 11 & 1 minute \\
Slide 12 & 1.5 minutes \\
Slide 13 & 1.5 minutes \\
\hline
\textbf{Total} & \textbf{$\sim$18 minutes} \\
\hline
\end{tabular}
\end{center}

\textit{Note: Adjust based on actual time limit}

\vspace{10pt}
\hrule
\vspace{10pt}

\section*{CONFIDENCE REMINDERS}

\begin{tipbox}
\begin{itemize}[leftmargin=*]
    \item[\checkmark] You know this material deeply---you wrote a comprehensive report on it
    \item[\checkmark] The PPT is well-structured and visually clear
    \item[\checkmark] Your work connects theory (Riemann) $\rightarrow$ method (SC) $\rightarrow$ applications (real world)
    \item[\checkmark] Professors want to see understanding, not perfection
    \item[\checkmark] Pauses are powerful---they show confidence
    \item[\checkmark] \textbf{You've got this!} \textbf{🎯}
\end{itemize}
\end{tipbox}

\vspace{20pt}
\begin{center}
\rule{0.8\textwidth}{0.4pt}\\[5pt]
\textbf{\Large END OF DELIVERY GUIDE}\\[5pt]
\textit{Best of luck with your presentation tomorrow!}
\end{center}

\end{document}
